\section{Cross-pollination and outlook}
\label{sec:cmp-outlook}

The two systems' gaps are not fundamental to characteristic two; each is
the price of a design choice the other made differently, and each has a
known lever.

\paragraph{Levers for \zinc{}.}
\begin{itemize}
  \item \emph{Proof size $\to$ WHIR.} The flat Brakedown proof is the
  cost of the linear-time encoder. A WHIR-style recursive proximity test
  \cite{whir} over the same code would make the proof $\polylog$,
  closing the structural proof-size gap independently of everything else.
  As a bonus it decouples proof size from the code rate, which would let
  \zinc{} raise the rate for a cheaper commit without paying in proof.
  \item \emph{Soundness $\to$ lookup adder.} The grand-product lookup
  adder (\Cref{sec:cmp-soundness}) closes Issue-1 and simultaneously removes
  the $12$ adder-carry Hadamards from the discharge --- one change that
  fixes soundness and trims the prover.
  \item \emph{Residual word-level discharge.} The big word-level win (the
  bit-axis univariate skip) is already in the oblong discharge; the
  residual is a row-axis eq-split worth $\sim$$5$--$8\%$ of the prover. It
  is protocol-visible (a fresh soundness argument), so research-grade, not
  a drop-in.
\end{itemize}

\paragraph{Levers for \binius{}.}
\begin{itemize}
  \item \emph{Verifier $\to$ holographic structure commitment.} The only
  way to make the $O(N)$ monster sweep succinct \emph{while keeping the
  general circuit} is a Spark/holographic commitment to the wiring
  \cite{spartan,lasso}, giving an $O(\log N)$ verifier. (A uniform AIR
  also does it, but only for uniform computations --- \Cref{sec:ana-air}.)
  A rayon parallelisation of the monster sweep is a measured but
  non-structural band-aid: it narrows the verify gap a constant factor
  without changing the $O(N)$ asymptotic.
  \item \emph{Prover $\to$ \zinc{}'s entry-wise shift.} Porting the free
  $\psi$ read-off of shifts would remove the Shift reduction that is
  $\sim$$49\%$ of \binius{}'s prover --- but doing so cleanly means
  adopting an $\Ftwo[X]$-style cell ring, i.e. converging toward \zinc{}.
\end{itemize}

\section{Conclusion}
\label{sec:conclusion}

\zinc{} and \binius{} agree on the hard decision --- prove SHA-256 in
characteristic two, refusing the prime-field witness blow-up --- and
disagree on three orthogonal downstream choices, each of which decides
one outcome. The uniform-AIR-versus-word-circuit choice decides the
verifier; the flat-versus-recursive proximity choice decides the proof;
the base-field-versus-$\mathrm{GF}(2^{128})$ reduction decides memory and
scaling. The prover, dominated by content both share, is a tie.
Measured on one machine: \zinc{} wins the verifier ($5.8\times$ to
$77\times$, structural and widening) and memory ($\sim$$2\times$);
\binius{} wins the proof ($3.8\times$ to $5.6\times$, structural and
widening) and is fully sound where \zinc{}'s adders are trusted; the
prover is a near-tie that tips to \zinc{} at scale.

Three observations resist the usual summaries and are worth carrying.
First, \binius{}'s linear verifier is caused by the \emph{cross-entry
wiring}, not the shifts: the entry-wise shift is succinct in both
systems, and the wiring is succinct only under a uniform arithmetization.
Second, the verifier and the prover are \emph{decoupled} in this family:
a succinct verifier costs the prover nothing to keep, and confers nothing
on it --- so neither ``slow the verifier to speed the prover'' (in either
system) nor ``an AIR makes \binius{}'s prover faster'' holds. Third, the
levers each system needs to close its gap point toward the other:
\zinc{}'s proof wants FRI-style recursion, and \binius{}'s prover wants
\zinc{}'s free shift --- the convergent endpoint is a binary-field STARK
with an $\Ftwo[X]$ cell ring, a uniform AIR, a base-field discharge, a
sound lookup adder, and a recursive proximity test.
